% \$ 2. Сечения многогранников плоскостями.

Сечения многогранника плоскостями могут иметь. сложную форму. На. рисунках
\ref{fig:8_2_11_a}-\ref{fig:8_2_11_c} приведены примеры сечений.
Видно, что они могут иметь неожиданную форму.

\begin{figure}\label{fig:8_2_11_a}
% стр 265 рис 11 a
\caption{Рис. 11 a) Сечение правильного тетраэдра плоскостью, параллельной ребрам
$AC$ и~$BD$ и~проходящей через середины остальных рёбер, представляет собой
квадрат. Докажите это.}
\end{figure}

% 266

\begin{figure}\label{fig:8_2_11_b}
% стр 266 рис 11 сдвоенный
\caption{Рис. 11 б) Сечение куба плоскостью, проходящей через его центр,
перпендикулярно главной диагонали, представляет собой правильный шестиугольник.
Докажите это.}
\end{figure}

\begin{figure}\label{fig:8_2_11_c}
% стр 266 рис 11 сдвоенный
\caption{Рис. 11 в) Сечение многогранника $\pi$ плоскостью , проходящей через рёбра
$AB$, $CD$ и~$EF$‚ представляет собой два прямоугольника (они заштрихованы),
соединённых отрезком $AB$. Отметим, что секущая плоскость $\alpha$ проходит
через внутренние точки многогранника (это внутренние точки заштрихованных
прямоугольников).}
\end{figure}

% 267

Сначала договоримся, что секущие плоскости должны проходить через внутренние
точки многогранника. Иначе мы должны относить к~числу сечений грани, ребра
и~вершины многогранника, что не отвечает привычному слову <<сечение>>.
Далее, совершенно необозримы сечения невыпуклых многогранников.
Это подтверждают примеры на рисунках
\ref{fig:8_2_11_a}-\ref{fig:8_2_11_c}.
%% рисунок 8_1_11 строенный
Поэтому мы ограничимся сечениями выпуклых многогранников.


\subsubsection{1.~Сечения плоскостями выпуклых многогранников.}

Напомним, что многогранник называется выпуклым, если он расположен по одну
сторону от плоскости каждой его грани.

Докажем следующую теорему.

\begin{Th}\label{th:8_2_1}
Теорема. Сечение выпуклого многогранника плоскостью, проходящей через его
внутреннюю точку, представляет собой выпуклый многоугольник
\end{Th}

Доказательство. Очевидно, сечение будет многоугольником. Нужно доказать,
что он выпуклый.

Пусть $\alpha$ --- секущая плоскость (рис.\ \ref{fig:8_2_12}) и $q$ --- многоугольник,
представляющий собой сечение нашего многогранника плоскостью $\alpha$.
Допустим, что $q$ --- невыпуклый многоугольник. Тогда он располагается в~плоскости
$\alpha$ по разные стороны от прямой, проходящей через некоторую его сторону $AB$
(вершины $C$ и~$D$‚ например, лежат в~плоскости $\alpha$ по разные стороны от
прямой $AB$).

Отрезок $AB$ лежит в~некоторой грани нашего многогранника.

\begin{figure}\label{fig:8_2_12}
% стр 297 рис 12
\caption{Рис. 12}
\end{figure}

% 268
Проведём. плоскость $\beta$ этой грани. Очевидно, что точки $C$ и~$D$ многогранника
лежат по разные стороны от плоскости $\beta$ ‚что противоречит условию выпуклости
многогранника. Таким образом, предположение о~невыпуклости многоугольника $q$
(сечения) ведёт к~противоречию с~условием выпуклости многогранника.
Следовательно, $q$ --- выпуклый многоугольник. Теорема доказана.

% 268

\subsubsection{2. Сечения правильного тетраэдра.}

В~начале этого параграфа мы говорили о~разнообразии форм плоских сечений
многогранников. На рисунках \ref{fig:8_2_11_a}-\ref{fig:8_2_11_c} указывалось,
что в~сечении правильного тетраэдра плоскостью может получиться квадрат.
Изучим более подробно форму сечений правильного тетраэдра плоскостями.

Отметим, что сечения тетраэдра плоскостями, проходящими через его внутренние
точки, представляют собой либо треугольники, либо четырехугольники (рис.\
\ref{fig:8_2_13}).

\begin{figure}\label{fig:8_2_13}
\caption{Рис.13. Сечения тетраэдра плоскостями либо треугольники,
либо четырехугольники.}
\end{figure}

Посмотрим, что представляют собой сечения правильного тетраэдра плоскостями,
параллельными двум его скрещивающимся рёбрам. Как мы только что отмечали,
одно из таких сечений представляет собой квадрат (рис.\ \ref{fig:8_2_11_a}).
Рассмотрим теперь сечение тетраэдра плоскостями, параллельными плоскости этого
квадрата. Для пояснения наших рассуждений мы обратимся к рисунку \ref{fig:8_2_14},
на котором изображен правильный тетраэдр (рис.\ \ref{fig:8_2_14_a}), ребро.
которого равно 1, и~развертка этого тетраэдра (рис.\ \ref{fig:8_2_14_b}).
Заметим, что если ребро правильного тетраэдра равно 1, периметр любого его
сечения плоскостями, параллельными двум скрещивающимся ребрам, будет равен 2.
Чтобы убедиться в этом, можно обратиться к~развертке. Видно, что граница
такого сечения --- прямоугольника, например, прямоугольника.$LMNO$, перейдет
в~отрезок $LL$ на развертке. Длина такого отрезка. равна периметру прямоугольника
$LMNO$ и равна 2.

% 269

% \begin{figure}\label{fig:8_2_14}

% TODO: использовать subfigure

% стр 269 рис 14 сдвоенный
% \caption{Рис.\ 14 Петля длины 2, через которую <<продевается>> тетраэдр}

\begin{figure}\label{fig:8_2_14_a}
\caption{Рис.14 а)~Сечения тетраэдра плоскостями, параллельными ребрам $AD$ и~$CB$
--- прямоугольники, имевшие периметр 2 (ребро тетраэдра равно 1).
Два таких прямоугольника $PQRS$ и~$LMNO$ изображены на рисунке. Границы этих
прямоугольников можно рассматривать как деформацию петли длины 2
через которую <<продлевается>>> тетраэдр.}
\end{figure}

\begin{figure}\label{fig:8_2_14_b}
\caption{Рис. 14 б)~Развертка тетраэдра с~ребром 1. Границы прямоугольников
 $PQRS$ и~$LMNO$ <<разворачиваются>> в~параллельные отрезки одинаковой длинны 2.}
\end{figure}

% \end{figure}

\begin{Note}
Замечание.Если представить себе тетраэдр сделанным из какого-либо материала,
например, из стали, и~поверхность его идеально отшлифованной, то, очевидно,
(если ребро равно 1) этот тетраэдр можно продеть сквозь гибкую и нерастяжимую
петлю, длина которой равна 2. Процесс <<продевания>> тетраэдра можно
представить себе так.
\end{Note}

% 210

Сначала через петлю, как бы сложенную из двух отрезков длины.2, продлевается
ребро $CB$, Затем эта петля, приняв форму прямоугольника (рис.\
\ref{fig:8_2_14_a}), скользит по поверхности тетраэдра и~затем снимается
с~него через ребро $AD$.

% 210

\subsubsection{3. Сечения куба.}

Сечения куба плоскостями, проходящими через его внутренние точки,  --- треугольники,
четырехугольники, пятиугольники, шестиугольники. Попытайтесь нарисовать
сечения такого вила.
На рисунке \ref{fig:8_2_11_b} показано сечение кубе плоскостью, проходящей
через его центр перпендикулярно главной диагонали. это сечение представляет
собой правильный шестиугольник (покажите это самостоятельно).

Изучим подробнее свойства этого сечения и~других сечений куба плоскостями,
параллельными плоскости указанного шестиугольника.

Для определённости будем считать, что ребро куба равно 2. Тогда сторона нашего
правильного шестиугольника-сечения будет равна (рис.\ \ref{fig:8_2_11_b} так,
что периметр правильного шестиугольника равен $6\sqrt{2}$  Оказывается,
если плоскость сечения перемещать параллельно плоскости нашего шестиугольника,
то до определенных пределов мы будем получать сечения, периметр которых также
равен $6\sqrt{2}$ чтобы убедиться в~этом, обратимся к~специальной развертке
куба (рис.\ \ref{fig:8_2_15_a}). В~текстах. под рисунками \ref{fig:8_2_15_a},
\ref{fig:8_2_15_b}, \ref{fig:8_2_15_c} указывается и~обосновывается
это интересное свойство куба. Отметим, что указанное свойство куба аналогично
свойству сечений правильного тетраэдра, установленного в~предыдущем пункте.

% 271

\begin{figure}\label{fig:8_2_15_a}
% стр 271 рис.15.
% Сечение плоскостью, параллельной плоскости правильного шестиугольника

\caption{Рис. 15 а) Куб с ребром 2. Секущая плоскость проходит через центр куба $O$
перпендикулярно главной диагонали $GD$.  В~сечении --- правильный
шестиугольник с~вершинами 1, 2, 3, 4, 5, 6. Периметр этого шестиугольника
равен $6\sqrt{2}$. Одно из сечений куба плоскостью, параллельной плоскости
шестиугольника, изображено штриховой линией.}
\end{figure}

% 272

\begin{figure}\label{fig:8_2_15_b}
% стр 272 рис 15_б

\caption{Рис. 15 б) Развёртка куба, изображенного на рисунке \ref{fig:8_2_15_a}.
Граница сечения --- правильного шестиугольника --- развертывается в~отрезок,
изображенный яркой линией. Вершины сечения --- занумерованные номерами
1, 2, 3, 4, 5, 6, 1 точки (точки 1, 1 --- отвечают одной и той же вершине 1).
Граница сечения плоскостью, параллельной плоскости правильного шестиугольника,
развертывается~отрезок, изображенный штрихом. Видно, что граница сечения
имеет периметр, равный длине этого отрезка, которая равна $6\sqrt{2}$.
Границы всех таких сечений (имеющих периметр $6\sqrt{2}$), заполняют на
развертке заштрихованный параллелограмм. При склеивании з развертки куба
стороны этого параллелограмма (обозначенные $BA$) склеются и~указанная
полоска расположится на кубе.Как это получится на самом деле,
изображено на рисунке \ref{fig:8_2_15_c}.}
\end{figure}

% 273

\begin{figure}\label{fig:8_2_15_c}
% стр 272 рис 15_в
\caption{Рис. 15 в)}
\end{figure}

Попытайтесь самостоятельно сформулировать и~обосновать аналогичные свойства
сечений ещё трех многогранников: правильного тетраэдра, додекаэдра и~икосаэдра.
На рисунках \ref{fig:8_2_16_a}, \ref{fig:8_2_16_b},
\ref{fig:8_2_17},
\ref{fig:8_2_18}
изображены эти многогранники и~их развёртки. Обратите внимание на указания
на рисунке \ref{fig:8_2_16_a}).

\begin{figure}\label{fig:8_2_16_a}
% стр 273 рис 16 сдвоенный
\caption{Рис. 16 а) Правильный октаэдр}
\end{figure}

\begin{figure}\label{fig:8_2_16_b}
% стр 273 рис 16 сдвоенный
\caption{Рис. 16 б) Развертка правильного октаэдра}
\end{figure}

% 274

\begin{figure}\label{fig:8_2_17}
% стр 274 рис 17 сдвоенный
\caption{Рис. 17 Правильный додекаэдр Развертка правильного додекаэдра
(масштаб уменьшен)}
\end{figure}

\begin{figure}\label{fig:8_2_18}
% стр 274 рис 18 сдвоенный
\caption{Рис. 18 Правильный икосаэдр Развертка правильного икосаэдра
(масштаб уменьшен)}
\end{figure}

Мы сейчас установим ещё одно свойство куба. Оказывается, в~данном кубе можно
прорезать такое отверстие, через которое может пройти куб, равный данному.
Для доказательства обратимся к~специальному сечению куба (рис.\ \ref{fig:8_2_19}).
Примем ребро куба $ABCDEFGH$  равным~1. Тогда сечение куба плоскостью,
проходящей через $EF$ и~$DC$ представляет собой прямоугольник со сторонами
$EF = 1$ и~$ED = \sqrt{2}$. Наглядно ясно, что при некотором <<шевелении>>
плоскости сечения можно получить квадрат, содержащийся в~кубе, со сторонами
больше 1. Поступим так: повернём плоскость прямоугольника  $EFCD$ так, чтобы
повернутая плоскость пересекалась с~передней гранью куба (грань $AEHD$)
по отрезку $QP$‚ параллельному $ED$, а~с~задней гранью (грань $BFGC$)
по отрезку $RS$ ‚параллельному $FC$.

% 275

\begin{figure}\label{fig:8_2_19}
% стр 275 рис 19
\caption{Рис. 19
$AD$ = 1, $AP = x$, $RG = x$. Если $x = \dfrac{3}{4}$, то $PQRS$ --- квадрат,
сторона $PQ$ которого равна $\dfrac{3\sqrt{2}}{4}$.
Так как это число больше 1, то в~квадрат $PQRS$ можно вписать квадрат со
стороной~1. Такой квадрат изображен штрихованной линией. Если в~данном кубе
прорезано отверстие, перпендикулярное плоскости квадрата $PQRS$ с~границей по
заштрихованной линии, то, очевидно, сквозь это отверстие пройдет куб,
равный данному.}
\end{figure}

Соединяя точки $P$, $Q$, $R$ и~$S$ отрезками, получим четырехугольник $PQRS$,
вершины которого расположены на ребрах куба и~поэтому сам четырехугольник
содержится в~кубе.
Убедимся, что можно так выбрать секущую плоскость, что четырехугольник
$PQRS$ будет квадрат. Наверное, проще всего поступить так. Введём три взаимно
перпендикулярных единичных вектора
$\overrightarrow{AD} = \vec{i}$,
$\overrightarrow{AB} = \vec{j}$,
$\overrightarrow{AE} = \vec{k}$
Кроме того, обозначим $AP = x$. Найдём теперь векторы $\overrightarrow{PQ}$
и~$\overrightarrow{PS}$ (векторы сторон четырехугольника $PQRS$).

% 276

Имеем

\begin{equation}
\left .
\begin{aligned}\label{eq:8_2_1}
\overrightarrow{PQ} &=
\overrightarrow{AQ} - \overrightarrow{AP} = x\vec{k} - x\vec{i} \\
\overrightarrow{PS} &=
\overrightarrow{PD} +
\overrightarrow{DC} +
\overrightarrow{CS} = (1 - x)\vec{i} + \vec{j} + (1 - x)\vec{k}
\end{aligned}
\quad \right\}
\end{equation}

Теперь вычислим длины этих векторов (длины сторон четырехугольника $PQRS$)
и~угол между ними. Так как длина вектора равна корню квадратному из скалярного
квадрата этого вектора,то из формул (\ref{eq:8_2_1}) получаем

\begin{equation}\label{eq:8_2_2}
PQ = x\sqrt{2}, \quad PS = \sqrt{1 + 2(1 - x)^{2}}.
\end{equation}

Для вычисления угла следует найти скалярное произведение
($\overrightarrow{PQ}$, $\overrightarrow{PS}$):

\begin{multline*}
(\overrightarrow{PQ}, \overrightarrow{PS}) =
(x\vec{k} - x\vec{i}, (1 - x)\vec{i} + \vec{j} + (1 - x)\vec{k}) = \\
x(1 - x) - x(1 - x) = 0.
\end{multline*}

Так как скалярное произведение ($\overrightarrow{PQ}$, $\overrightarrow{PS}$)
равно нулю, то векторы $\overrightarrow{PQ}$ и~$\overrightarrow{PS}$
перпендикулярны и,~следовательно, четырехугольник $PQRS$ --- прямоугольник
(при любом выборе $AP = x$). Если стороны этого прямоугольника будут равны,
он превратится в~квадрат. Найдем значение $x$ ‚при котором $PQRS$ --- квадрат.
Из (\ref{eq:8_2_2}) получаем нужное соотношение для $x$ (приравниваем $PQ$
и $PS$):

\begin{equation*}
2x^{2} = 1 + 2(1 - x)^{2}.
\end{equation*}

Отсюда

\begin{equation}\label{eq:8_2_3}
AQ = AP = x = \dfrac{3}{4}.
\end{equation}

При этом значении $x$ четырехугольник $PQRS$ --- квадрат.
Найдем его $PQ$. Из треугольника $APQ$ имеем

\begin{equation*}
PQ = \sqrt{AP^{2} + AQ^{2}}.
\end{equation*}

Отсюда и~из соотношения (\ref{eq:8_2_3}) получаем

\begin{equation}\label{eq:8_2_4}
PQ = \sqrt{\dfrac{9}{8}} = \dfrac{3\sqrt{2}}{4}.
\end{equation}

Из полученного выражения для $PQ$ видно, что $PQ > 1^{x}$\footnote{
Так как $\sqrt{2} = 1,41\dots$  то из \ref{eq:8_2_4} следует, что приближенное
значение $PQ$ равно 1,056.}
Можно сделать следующий вывод: так как длина стороны квадрата $PQRS$ больше 1,
то в~этом квадрате целиком поместится квадрат со стороной 1 (и~даже чуть больше,
чем 1). На рисунке \ref{fig:8_2_19} такой квадрат изображен штрихованной линией.
Прорежем теперь в~данном кубе отверстие в~направлении, перпендикулярном плоскости
квадрата $PQRS$ и~с~границей заштрихованной линии (рис.\ \ref{fig:8_2_19}),
то через это отверстие пройдет куб, равный данному.

\begin{Note}
Из наших рассуждений следует, что в~данном кубе можно вырезать отверстие, через
которое пройдёт куб, длина ребра которого больше, чем у данного.
\end{Note}
